
\bta{2008}

\begin{enumerate}
	%\renewcommand{\labelenumi}{\arabic{enumi}.}
	% A(\Alph) a(\alph) I(\Roman) i(\roman) 1(\arabic)
	%设定全局标号series=example	%引用全局变量resume=example
	%[topsep=-0.3em,parsep=-0.3em,itemsep=-0.3em,partopsep=-0.3em]
	%可使用leftmargin调整列表环境左边的空白长度 [leftmargin=0em]
	\item
（$30'$）写成氢原子的束缚态能级，所有量子数以及取值范围，求出其简并度。


\banswer{
	忽略精细结构时，氢原子的束缚态能量为
	\[
		\boxed{E_n=-\frac{\mu e^4}{2(4\pi\varepsilon_0)^2\hbar^2n^2}},
		\qquad n=1,2,3,\ldots,
	\]
	其中 $\mu$ 为约化质量。空间波函数的量子数范围为
	\[
		l=0,1,\ldots,n-1,\qquad m=-l,-l+1,\ldots,l.
	\]
	对固定 $n$，空间简并度为
	\[
		g_n^{(\mathrm{space})}=\sum_{l=0}^{n-1}(2l+1)=n^2.
	\]
	若同时计入电子自旋量子数 $s=1/2$、$m_s=\pm1/2$，而忽略自旋相关相互作用，则总简并度为
	\[
		\boxed{g_n=2n^2}.
	\]
}


\item 
（$30'$）一个粒子质量为$\mu$，在一势能环中运动，势能
$$V=\begin{cases}0,&0<\varphi<\varphi_0\\ \infty,&other\end{cases}$$
求粒子运动的本征值和本征函数。


\banswer{
	设圆环半径为 $R$，转动惯量为 $I=\mu R^2$。在允许的角区间内，定态方程为
	\[
		-\frac{\hbar^2}{2I}\frac{d^2\psi}{d\varphi^2}=E\psi,
		\qquad \psi(0)=\psi(\varphi_0)=0.
	\]
	写通解 $\psi=A\sin(k\varphi)+B\cos(k\varphi)$，其中
	$k=\sqrt{2IE}/\hbar$。左端条件使 $B=0$，右端使 $k\varphi_0=n\pi$。
	由 $|A|^2\int_0^{\varphi_0}\sin^2(n\pi\varphi/\varphi_0)d\varphi=1$ 得
	$|A|=\sqrt{2/\varphi_0}$。因此
	\[
		\boxed{E_n=\frac{\hbar^2\pi^2n^2}{2I\varphi_0^2}},
		\qquad n=1,2,3,\ldots,
	\]
	相对于角度测度 $d\varphi$ 归一化的本征函数为
	\[
		\boxed{\psi_n(\varphi)=
		\begin{cases}
		\sqrt{\dfrac{2}{\varphi_0}}\sin\dfrac{n\pi\varphi}{\varphi_0},&0<\varphi<\varphi_0,\\[6pt]
		0,&\text{其他角区间}.
		\end{cases}}
	\]
	题干未给出圆环半径，故能量只能用 $I=\mu R^2$ 表示。
}


\item 
（$30'$）求在$H=\begin{pmatrix}1&\lambda &0\\ \lambda &3&0 \\ 0&0& \lambda-2\end{pmatrix}$中粒子的本征值，设$\lambda \ll 1$，利用微扰求其本征值（精确到二级近似）并与精确求解相比较。


\banswer{
	矩阵的精确本征值为
	\[
		E_1=2-\sqrt{1+\lambda^2},\qquad
		E_2=2+\sqrt{1+\lambda^2},\qquad
		E_3=\lambda-2.
	\]
	当 $|\lambda|\ll1$ 时展开到二级，
	\[
		\boxed{E_1=1-\frac{\lambda^2}{2}+O(\lambda^4),\qquad
		E_2=3+\frac{\lambda^2}{2}+O(\lambda^4),\qquad
		E_3=-2+\lambda}.
	\]
	用微扰论也得到同样结果：$\lambda=0$ 时三个能级为 $1,3,-2$；微扰只以矩阵元 $\lambda$ 耦合前两个态，故它们的一级修正为零、二级修正分别为
	$\lambda^2/(1-3)=-\lambda^2/2$ 与 $\lambda^2/(3-1)=\lambda^2/2$；第三态仅有对角一级修正 $\lambda$。
}




\newpage
\item 
（$30'$）两个自旋为$\frac{1}{2}\hbar$的粒子，两个粒子分别为$\chi_1=\begin{pmatrix}1\\0\end{pmatrix}$，$\chi_1=\begin{pmatrix}\cos\theta e^{-i\omega t}\\ \sin\theta e^{-i\omega t}\end{pmatrix}$，求系统处于单态和三态的概率。

\banswer{
	把第二个粒子的状态记为
	\[
		|\chi_2\rangle=e^{-i\omega t}
		(\cos\theta|\uparrow\rangle+\sin\theta|\downarrow\rangle).
	\]
	总态为 $|\uparrow\rangle_1|\chi_2\rangle$。投影到单态
	$|0,0\rangle=(|\uparrow\downarrow\rangle-|\downarrow\uparrow\rangle)/\sqrt2$，得
	\[
		\boxed{P_{S=0}=\frac12\sin^2\theta}.
	\]
	三重态总概率为
	\[
		\boxed{P_{S=1}=1-P_{S=0}=1-\frac12\sin^2\theta}.
	\]
	更具体地，$P_{1,1}=\cos^2\theta$、$P_{1,0}=\sin^2\theta/2$、$P_{1,-1}=0$。题给的共同相位 $e^{-i\omega t}$ 不影响这些概率。
}


\item 
（$30'$）处在一维谐振子势基态的粒子受到微扰$H'=\lambda x\delta(t_0)$作用，求跃迁到其他各激发态的总概率和仍处在基态的概率。已知：
$xH_n=\frac{1}{\sqrt a}\left[\sqrt{\frac{n}{2}}H_{n-1}+\sqrt{\frac{n+1}{2}}H_{n+1}\right]$。

\banswer{
	跨过 $t=t_0$ 的瞬时微扰时，波函数受到冲量演化算符
	\[
		U_{\mathrm{kick}}=\exp\!\left(-\frac{i\lambda x}{\hbar}\right)
	\]
	作用；题面 $\delta(t_0)$ 按在 $t=t_0$ 的完整脉冲 $\delta(t-t_0)$ 理解。
	写 $x=x_0(a+a^\dagger)$、$x_0=\sqrt{\hbar/(2m\omega)}$，
	并令 $\alpha=-i\lambda x_0/\hbar$。因 $\alpha$ 纯虚，
	\[
	 U_{\rm kick}=e^{\alpha a^\dagger-\alpha^*a}
	 =e^{-|\alpha|^2/2}e^{\alpha a^\dagger}e^{-\alpha^*a}.
	\]
	这里用了两指数的对易子为常数的 BCH 公式。又 $a|0\rangle=0$，故
	\[
	 U_{\rm kick}|0\rangle=e^{-|\alpha|^2/2}
	 \sum_{n=0}^\infty\frac{\alpha^n}{\sqrt{n!}}|n\rangle.
	\]
	这是相干态，参数为
	\[
		\alpha=-\frac{i\lambda}{\sqrt{2m\hbar\omega}},
		\qquad |\alpha|^2=\frac{\lambda^2}{2m\hbar\omega}.
	\]
	故跃迁到第 $n$ 个能级的概率为 Poisson 分布
	\[
		P_n=e^{-|\alpha|^2}\frac{|\alpha|^{2n}}{n!}.
	\]
	特别地，仍处于基态和跃迁到所有激发态的总概率分别为
	\[
		\boxed{P_0=\exp\!\left(-\frac{\lambda^2}{2m\hbar\omega}\right)},
		\qquad
		\boxed{P_{\mathrm{exc}}=1-P_0}.
	\]
	在弱冲量近似下，$P_{\mathrm{exc}}\simeq\lambda^2/(2m\hbar\omega)$，且最低阶只有第一激发态有贡献。
}



\end{enumerate}
